\section{Diagnostic revision: evidence and interpretation}
This supplement places the later diagnostic methods and explanatory figures before the complete historical tables. The neural experiments and the original-grid Ridge evidence are preserved. New regularization fits are post hoc sensitivities on already observed performance splits, not new prespecified confirmation. File-based tests below concern numerical and bookkeeping consistency, not a count of independent experiments.

\subsection{Yulara actual-matrix diagnosis}
The saved Train-fitted processor, design means/scales, coefficients, target transformation, and timestamps reconstruct the original Test arrays. The objective is $\|Y-XB\|_F^2+\alpha\|B\|_F^2$ after standardization and target centering, with an unpenalized intercept. A Cholesky solve uses the same actual 69333-by-504 or 69333-by-792 matrix. For the original selected parameters 0.1 and 1000, the maximum Test difference from the saved spectral solution is below $9\times10^{-9}$ kW. Normal-equation relative residuals and objective comparisons are supplied, together with the full per-feature Train scales and later offsets.

The origin-position weather-missing flags occur in four eligible Train windows. Their small scale, 0.007595, maps an observed missing flag to about 132 standardized units. Validation exposure is 18.56\% and Test exposure 5.50\%. The two weather masks coincide and their prediction components are correlated; their separate RMS components must not be added as independent MSE contributions. These linear components and the concentration of errors in a small group of windows support a missing-pattern extrapolation explanation. No window is removed and no prediction is clipped. Diagnostic quantiles use every finite target in saved candidate trajectories; headline twelve-hour metrics use the stricter complete-H144 and Daily-finite intersection.

Seven original alpha choices lie at grid endpoints. After diagnosing the actual matrices, a common $10^{-4},\ldots,10^8$ grid was fixed for all five systems and A/B variants. The same original Train transformations and full-H144 Validation objective select each parameter. Every Validation candidate is retained. The expanded choices and Test errors are supplementary sensitivity results, with the original grid still reported. Detailed solver records distinguish original-model verification from low-alpha numerical checks in the supplemental fits.

\subsection{Qcells exclusion attribution}
Every five-minute coordinate within a split is a candidate origin. Overlapping flags identify insufficient input history or target tail, future missing/nonfinite labels, finite negative future labels, invalid origin power and unavailable Daily points. Exclusive accounting gives boundary precedence, then future nonfinite, future negative, origin invalid and retention. Daily is a separate point intersection; it is not used to explain primary origin deletion. Tables distinguish origin counts, origin--lead target pairs and unique physical target timestamps.

For Test one-hour support, complete-twelve-hour conditioning removes 3467 origins: 117 boundary cases and 3350 future-negative cases under the exclusive order. The negative flag overlaps 65 of those boundary cases, giving 3415 flagged origins overall. There are no future nonfinite-label exclusions. Fifty-four raw negative records (52 at hour 18, two at hour 7) affect many overlapping long trajectories. Their values range from approximately $-0.010933$ kW to numerical values near zero; no physical or invalid-code interpretation is inferred. The frozen rule remains unchanged. Train and Validation exhibit the same selection process. Original neural Train/Validation label-completeness counts can differ from evaluation because the historical Alice fitting routine does not require a valid origin power.

\subsection{Ridge paired temporal effects}
All comparisons use complete H144 origins and an identical Daily-finite target intersection across A, B, Daily and Inverted seeds 42/43/44. Full, power-active and low-power use the same unchanged Train threshold. Blocks are anchored at the Test split's midnight in the frozen coordinate and assigned by forecast origin, not target time. Fixed EST is retained at NIST; Yulara has provider-local coordinates. Empty blocks remain in the resampling universe. The 48-hour main scheme and 24/72-hour sensitivity use 2000 draws and random seed 20260908.

Each draw sums block SSE and valid counts, then computes RMSE. B-minus-reference is negative when B is better. Each real neural seed is compared separately; the mean contrast averages per-seed effects under the same draw. Deterministic Ridge has no artificial seed repetitions. The 48-hour Test series contains 12 Alice blocks or 46 external blocks; some power subsets have fewer nonempty blocks. Full block counts, SSE, actual target-index support and effects are in the lightweight evidence. Overlapping trajectories can still couple adjacent blocks, and intervals condition on fitted predictions and observed dates. They do not include retraining or guarantee unobserved-year performance.

\subsection{Actual versus historical validation}
This diagnostic revision checks saved Ridge predictions against reconstructed matrices. It compares solver objectives and residuals, exact timestamp alignment, Qcells exclusion accounting and paired block arithmetic. SHA-256 inventories compare the original data, frozen results, checkpoints and original prediction arrays before and after. The old M1-R/M2 full suites and original neural checkpoint reproduction remain historical records, not new claims. Memory chunking and datetime-resolution normalization in the diagnostic code do not alter the fixed EST or provider-local coordinate. No neural training occurs.

At the expanded NIST A choice of $10^{-4}$, an independent augmented QR solve has full column rank 504. Its maximum prediction difference from the spectral solution is 0.002390 kW; the full-scope RMSE difference is $-1.835\times10^{-7}$ kW. This small-alpha numerical sensitivity is retained in the diagnostic evidence, rather than replacing either prediction array. It does not affect the original-grid NIST B contrasts reported in the main text.


\clearpage\section{Post hoc review: support, uncertainty and linear references}
These analyses were planned after the original performance results were known; they do not change the frozen neural experiments. H12/H48/H96/H144 are cumulative 1/4/8/12 h windows. Common-H144 comparisons retain only complete twelve-hour origins and evaluate every prefix there. Fixed-lead curves score individual leads on that same support. Daily comparisons use the same elementwise target intersection for every method. The tables in the accompanying review CSVs include both origin and target counts for every support.

Power-active targets have true power above 1\% of Train maximum; low-power is the complement, not night. In that preceding review, all 2400 tested full/active/low-power squared-error identities passed on saved neural and reference arrays. Full MSE is the sum of subset SSE divided by the full target count. The per-seed decomposition file also records MAE, bias, counts and each weighted MSE contribution. Its neural summaries must average per-seed metrics rather than square an averaged RMSE.

The paired block procedure uses fixed 24/48/72 h nonoverlapping origin blocks anchored at Test midnight, 2000 resamples, and random state 20260908. Forty-eight hours is the primary block length. Empty blocks are retained. Methods and seeds share each draw. Per-seed RMSE is recomputed from summed SSE and count; the reported statistic is the mean of per-seed reference skill. Percentile intervals condition on the fitted models and observed dates. Dependence can persist across block boundaries, and Alice has much shorter temporal support than the external sites. No p-values or independent-comparison interpretation are attached to win counts.

Ridge A uses 72 times the neural input width in flattened, Train-standardized features (1224 at Alice and 504 externally). Ridge B adds 144 lagged power values at target time minus 24 h and 144 missingness indicators, giving 1512 or 792 features. Train medians fill missing lag inputs. The same Train-only numerical imputer, masks, Isolation Forest and neural feature transformations precede these linear designs. A separate standardizer is fitted to each final design. The target transformation is fitted to valid Train power. A single regularization parameter per system and variant is selected from 0.1, 1, 10, 100, 1000 on complete H144 Validation global MSE. Deterministic spectral ridge solutions include an unpenalized intercept; no Test values choose regularization. All ten fitted references and all 50 Validation candidate scores are retained, including the unfavorable Yulara results.

Monthly summaries assign complete H144 origins to October, November or December. Fixed cases use the first eligible origin at or after the 15th at 10:00. Additional best/worst cases explicitly use post hoc seed42 Daily-error differences on complete Daily-valid trajectories. Each displayed twelve-hour curve belongs to one origin. Neither weather conditions nor operation events are inferred from appearance.

The preceding September 8 review re-read all 60 training histories and verified each recorded best Validation epoch. It independently checked 480 Ridge metric rows, all ten selected parameters and finite prediction shapes, and five synthetic spectral/direct solver equivalences: 565 checks passed, zero failed or skipped. The 43/17/10 historical protocol/artifact suites and 10432 M2 evidence comparisons reported above are historical audit results, not newly repeated full suites in this review. No neural fit, checkpoint update or prediction regeneration was performed. A further 6720 metric/count comparisons between the new point-array reductions and the unchanged historical CSVs passed, with zero failures or skips. Eight additional synthetic forward checks confirmed both input widths and participation of parameter-owning modules, without loading checkpoints. The saved-prediction audits and their explicit scope accompany the review materials.

